\subsection{Weyl Bounds and Rank-One Spectral Shifts}

\begin{theorem} \label{thm:weyl-eigenvalue-inequalities-hermitian-sums} \hlt{Weyl Eigenvalue Inequalities}\\
Let $A,B\in\operatorname{Mat}(n,n,\mathbb{C})$ be Hermitian. For $1\leq i\leq n$,
\begin{align*}
\lambda_i(A+B) &\leq \lambda_{i+j}(A)+\lambda_{n-j}(B),
& 0 &\leq j\leq n-i,\\
\lambda_{i-j+1}(A)+\lambda_j(B) &\leq \lambda_i(A+B),
& 1 &\leq j\leq i.
\end{align*}
\begin{proof}
Let $\mathbf{u}_k,\mathbf{v}_k,\mathbf{w}_k$ be increasing eigenbases for $A,A+B,B$.\\
Fix $i$ and $0\leq j\leq n-i$, and set $\mathcal{U}=\operatorname{span}(\mathbf{u}_1,\ldots,\mathbf{u}_{i+j})$, $\mathcal{V}=\operatorname{span}(\mathbf{v}_i,\ldots,\mathbf{v}_n)$, $\mathcal{W}=\operatorname{span}(\mathbf{w}_1,\ldots,\mathbf{w}_{n-j})$.\\
Their dimensions sum to $2n+1$, so choose a unit $\mathbf{x}\in\mathcal{U}\cap\mathcal{V}\cap\mathcal{W}$.\\
Then $\lambda_i(A+B)\leq \mathbf{x}^*(A+B)\mathbf{x} =\mathbf{x}^*A\mathbf{x}+\mathbf{x}^*B\mathbf{x} \leq \lambda_{i+j}(A)+\lambda_{n-j}(B)$.\\
The lower bound follows by applying this upper bound to $-A$ and $-B$ and using $\lambda_k(-C)=-\lambda_{n-k+1}(C)$.\\
\hlt{Structural Commentary}: Weyl bounds come from intersecting the spectral flags of $A$, $B$, and $A+B$.\\
\hlt{Key Insight}: One shared test vector controls all three Rayleigh quotients.
\end{proof}
\end{theorem}

\begin{method} \label{meth:weyl-inequality-method} \hlt{Weyl Bound Method}\\
To bound eigenvalues of $A+B$:
\begin{enumerate}[label=(\roman*)]
\setlength{\itemsep}{0pt}
\item choose the target index $i$;
\item use the upper Weyl family with $j=0,\ldots,n-i$ to get upper bounds;
\item use the lower Weyl family with $j=1,\ldots,i$ to get lower bounds;
\item for positive semidefinite or low-rank $B$, insert the known zero or nonnegative eigenvalues to recover monotonicity or interlacing.
\end{enumerate}
\end{method}

\begin{theorem} \label{thm:brauer-rank-one-eigenvalue-shift} \hlt{Brauer Rank-One Eigenvalue Shift}\\
Let $A\in\operatorname{Mat}(n,n,\mathbb{C})$, let $A\mathbf{x}=\lambda\mathbf{x}$ with $\mathbf{x}\neq\mathbf{0}$, and let $\mathbf{y}\in\mathbb{C}^n$. If $\lambda,\lambda_2,\ldots,\lambda_n$ are the eigenvalues of $A$ with multiplicity, then the eigenvalues of $A+\mathbf{x}\mathbf{y}^*$ are $\lambda+\mathbf{y}^*\mathbf{x},\lambda_2,\ldots,\lambda_n$.\\
Moreover, $\mathbf{x}$ is an eigenvector of $A+\mathbf{x}\mathbf{y}^*$ corresponding to $\lambda+\mathbf{y}^*\mathbf{x}$.
\begin{proof}
Put $\mathbf{u}=\mathbf{x}/\norm{\mathbf{x}}_2$. By Schur triangularisation, choose a unitary $U=[\mathbf{u}\ U_2]$ such that $U^*AU=\begin{bmatrix}\lambda&*\\0&T'\end{bmatrix}$, where the diagonal entries of $T'$ are $\lambda_2,\ldots,\lambda_n$. Also $U^*(\mathbf{x}\mathbf{y}^*)U
=
\begin{bmatrix}
\mathbf{y}^*\mathbf{x}&*\\
0&0
\end{bmatrix}$.
Therefore $U^*(A+\mathbf{x}\mathbf{y}^*)U$ is upper triangular with diagonal entries $\lambda+\mathbf{y}^*\mathbf{x},\lambda_2,\ldots,\lambda_n$. Finally, $(A+\mathbf{x}\mathbf{y}^*)\mathbf{x}=A\mathbf{x}+\mathbf{x}(\mathbf{y}^*\mathbf{x})=(\lambda+\mathbf{y}^*\mathbf{x})\mathbf{x}$.
\hlt{Structural Commentary}: A rank-one eigenvector update changes only the first Schur diagonal entry.\\
\hlt{Key Insight}: Align the update with the first coordinate; only one diagonal term shifts.
\end{proof}
\end{theorem}

\begin{figure}[H]
  \centering
  \begin{tikzpicture}[
  >=Stealth,
  every node/.style={font=\small},
  mat/.style={draw=gray!55, fill=gray!5, rounded corners=2pt, align=center, inner sep=6pt},
  shiftbox/.style={draw=blue!60!black, fill=blue!7, rounded corners=2pt, align=center, inner sep=6pt},
  arr/.style={->, line width=0.75pt, blue!70!black}
]
  \node[mat] (schur) at (0,0)
    {$U^*AU=\begin{bmatrix}\lambda&*\\0&T'\end{bmatrix}$};
  \node[shiftbox] (rank) at (4,0)
    {$U^*\mathbf{x}\mathbf{y}^*U=\begin{bmatrix}\mathbf{y}^*\mathbf{x}&*\\0&0\end{bmatrix}$};
  \node[mat] (sum) at (8.2,0)
    {$\begin{bmatrix}\lambda+\mathbf{y}^*\mathbf{x}&*\\0&T'\end{bmatrix}$};

  \draw[arr] (schur) -- (rank);
  \draw[arr] (rank) -- (sum);
\end{tikzpicture}

  \caption{Brauer's update shifts one Schur diagonal entry.}
  \label{fig:brauer-rank-one-shift}
\end{figure}

\begin{corollary} \label{cor:damped-rank-one-eigenvalue-reset} \hlt{Damped Rank-One Eigenvalue Reset}\\
Let $A\in\operatorname{Mat}(n,n,\mathbb{C})$ have eigenpair $(\lambda,\mathbf{x})$, let $\mathbf{y}^*\mathbf{x}=1$, and let $\tau\in\mathbb{C}$. If $\lambda,\lambda_2,\ldots,\lambda_n$ are the eigenvalues of $A$, then $A_{\tau}\defeq \tau A+(1-\tau)\lambda\mathbf{x}\mathbf{y}^*$ has eigenvalues $\lambda,\tau\lambda_2,\ldots,\tau\lambda_n$.
\begin{proof}
The eigenvalues of $\tau A$ are $\tau\lambda,\tau\lambda_2,\ldots,\tau\lambda_n$. Since $A_{\tau}=\tau A+\mathbf{x}((1-\overline{\tau})\overline{\lambda}\mathbf{y})^*$, Theorem \ref{thm:brauer-rank-one-eigenvalue-shift} shifts only $\tau\lambda$ by $(1-\tau)\lambda\mathbf{y}^*\mathbf{x}=(1-\tau)\lambda$. Thus $\tau\lambda$ becomes $\lambda$, while the other eigenvalues stay $\tau\lambda_i$.\\
\hlt{Structural Commentary}: Damping scales the spectrum; Brauer's update restores one eigenvalue.\\
\hlt{Key Insight}: Brauer's theorem changes one spectral coordinate without moving the rest.
\end{proof}
\end{corollary}
